\documentclass[10pt,a4paper]{article}
\usepackage[margin=1.45cm]{geometry}
\usepackage{graphicx,booktabs,tabularx,array,float,enumitem,hyperref,caption,microtype,titlesec}
\setlength{\parindent}{0pt}
\setlength{\parskip}{2.2pt}
\setlist[itemize]{leftmargin=*,noitemsep,topsep=1pt}
\captionsetup{font=footnotesize,skip=2pt}
\hypersetup{colorlinks=true,linkcolor=black,urlcolor=blue,citecolor=black}
\titlespacing*{\section}{0pt}{4pt}{1.5pt}
\titleformat{\section}{\normalsize\bfseries}{\thesection.}{0.4em}{}
\newcolumntype{Y}{>{\raggedright\arraybackslash}X}

\begin{document}
\begin{center}
{\large\textbf{MIT 805: Big Data Semester Project}}\\[-1pt]
{\normalsize\textbf{Part 1 -- Data Collection \& Analysis: Amazon Reviews 2023}}\\[2pt]
\textbf{Group 6:} RY (Rabelani) Muthadzwi \quad | \quad M (Makungu) Ndlovu
\end{center}
\vspace{-2mm}\hrule\vspace{1.5mm}

\section{Dataset Provenance and Description}
We analyse the \textbf{Clothing, Shoes and Jewelry} review file from McAuley Lab's Amazon Reviews 2023 collection \cite{hou2024}. The public repository provides review text, title, rating, product and pseudonymous user identifiers, timestamp, helpful-vote count, image references, and verified-purchase status \cite{amazonhf}. The hosted file is JSON Lines (JSONL), an appropriate record-oriented format for distributed ingestion. The collection covers Amazon activity from May 1996 to September 2023; our processed prefix spans December 1999 to March 2023.

The maintainers state that the data are provided primarily for research but do not assign a standalone licence \cite{terms}. They explicitly permit academic thesis use, subject to ethical and legal responsibilities. We therefore use the data only for non-commercial coursework, cite the creators, do not redistribute review-level data, and do not attempt to re-identify pseudonymous reviewers.

\section{Dataset Size and Characteristics}
The selected raw file was verified through the Hugging Face API and downloaded in full. A line-safe prefix was then produced: complete JSONL records were copied until the processing file exceeded 3~GiB, preventing a truncated final record.

\begin{table}[H]
\centering\footnotesize
\caption{Measured dataset scale and characteristics.}
\begin{tabularx}{\textwidth}{@{}p{2.8cm}Y@{}}
\toprule
\textbf{Characteristic} & \textbf{Measured value} \\
\midrule
Raw dataset & 27,810,080,533 bytes (27.810 GB; 25.900 GiB), satisfying the required 25--40 GB range. \\
Working dataset & 27,810,080,533 bytes: the complete downloaded category file, exceeding 12 GB. \\
Processing dataset & 3,221,225,482 bytes (3.221 GB; 3.000 GiB); 7,231,543 records processed with PySpark 4.0.0. \\
Structure & JSONL; 10 source variables and 4 derived variables (review time/year and text/title lengths). \\
Observed period & 1 December 1999--25 March 2023 in the processing subset. \\
\bottomrule
\end{tabularx}
\end{table}

\section{Data Quality and Preparation}
PySpark inferred the JSON schema; ratings, votes, and timestamps were explicitly cast, millisecond timestamps converted to calendar time, and text/title lengths derived. All 7,231,543 rows were complete across the 14 analysed fields, and ratings were within 1--5 with no negative helpful-vote counts. However, 495 reviews (0.0068\%) had empty text, 18 exceeded 10,000 characters, and helpful votes were strongly right-skewed (median 0, mean 0.919, maximum 41,199). Grouping by user, parent product, timestamp, and text identified 23,544 potential duplicate rows (0.3256\%). These are flagged rather than automatically removed because Amazon may merge reviews for product variants.

Quality limitations remain despite field completeness. Reviews are voluntary, so purchasers who post are unlikely to represent all customers. Ratings can also reflect selection effects, platform design, incentives, or manipulation. Finally, a file prefix is computationally reproducible but may differ temporally or compositionally from the complete category; processed-subset estimates must not be presented as exact full-file totals.

\section{Exploratory Data Analysis}
Ratings were highly positive: 4,647,623 five-star reviews (64.27\%) and 1,085,590 four-star reviews (15.01\%) produced a mean of 4.253 and median of 5. Only 6.53\% were one-star. Review text averaged 171 characters (median 104), showing that a small number of very long reviews inflate the mean.

Annual volume increased from 106,455 reviews in 2013 to a peak of 1,133,095 in 2021, then reached 1,054,750 in 2022. The 2023 count is partial and must not be interpreted as a full-year decline. Verified purchases comprised 6,688,231 records (92.48\%) and had a slightly higher mean rating than unverified reviews (4.255 versus 4.228); the small difference is descriptive, not causal. Product-level concentration was also substantial: the most-reviewed parent ASIN had 37,409 reviews, over six times the second-ranked product. Supporting plots appear in Appendix~\ref{app:eda}.

\section{The 7 Vs of Big Data}
\begin{table}[H]
\centering\scriptsize
\caption{Evidence-based application of the 7 Vs.}
\begin{tabularx}{\textwidth}{@{}p{1.3cm}Y@{}}
\toprule
\textbf{V} & \textbf{Dataset evidence and importance} \\
\midrule
Volume & A 27.810 GB working file and 3.221 GB processing subset yielded 7.23 million rows, motivating distributed ingestion and aggregation. \\
Velocity & Reviews accumulate as customer events; yearly counts grew beyond one million in 2019, 2020, 2021, and 2022, creating temporal drift. \\
Variety & Text, numeric ratings/votes, Booleans, timestamps, arrays, and identifiers require mixed parsing and transformations. \\
Veracity & Empty text, 23,544 potential duplicates, extreme vote/text values, self-selection, and possible manipulation can bias conclusions. \\
Value & Rating, demand, verification, and helpfulness patterns support quality monitoring, merchandising, and customer-experience decisions. \\
Variability & Annual mean ratings ranged from 4.154 in 2010 to 4.311 in 2019 after volumes became substantial; products also varied in rating and review count. \\
Visualization & Rating and time-series charts compress millions of records into interpretable distributions and trends without moving raw data to Pandas. \\
\bottomrule
\end{tabularx}
\end{table}

\section{Business/Societal Value and Limitations}
Retailers and manufacturers can use the results to identify unusually low-rated, high-volume products; prioritise recurring quality concerns; and monitor changes in customer experience. Review helpfulness and verification may support ranking and trust systems, while temporal patterns can guide capacity planning. These applications require caution: association is not causation, reviewers are self-selected, pseudonymous records cannot establish demographic representativeness, and the line-safe processing prefix may not mirror the full file. Ethical use therefore requires aggregate reporting, privacy protection, and validation before operational decisions.

\textbf{Reproducibility:} notebook, dependencies, and data-access instructions are available at \url{https://github.com/Rabelani-byte/MIT805-Big-Data-Project}. Raw data are excluded from GitHub.

\begin{thebibliography}{9}\footnotesize
\bibitem{hou2024} Hou et al. (2024), \textit{Bridging Language and Items for Retrieval and Recommendation}. \url{https://arxiv.org/abs/2403.03952}
\bibitem{amazonhf} McAuley Lab, \textit{Amazon Reviews 2023 -- Hugging Face}. \url{https://huggingface.co/datasets/McAuley-Lab/Amazon-Reviews-2023}
\bibitem{terms} McAuley Lab, \textit{Dataset licence discussion}. \url{https://huggingface.co/datasets/McAuley-Lab/Amazon-Reviews-2023/discussions/1}
\end{thebibliography}

\clearpage
\appendix
\section{Supporting EDA}\label{app:eda}
\begin{figure}[H]
\centering
\includegraphics[width=.78\textwidth]{rating_distribution.png}
\caption{Rating distribution for the 7,231,543-record processing subset.}
\end{figure}
\begin{figure}[H]
\centering
\includegraphics[width=.82\textwidth]{review_volume_by_year.png}
\caption{Annual review volume. The 2023 observation is a partial year.}
\end{figure}
\end{document}

